\documentclass[12pt]{article}
\usepackage[utf8]{inputenc}
\usepackage[T1]{fontenc}
\usepackage[margin=1in]{geometry}
\usepackage{setspace}
\usepackage{xcolor}
\usepackage[colorlinks=true, urlcolor=blue!60!black, linkcolor=blue!60!black]{hyperref}
\setlength{\parindent}{2em}
\setlength{\parskip}{1.25em}
\renewcommand{\baselinestretch}{1.0}
\begin{document}
\begin{center}
{\LARGE \textit{Prigg v. Pennsylvania}, 41 U.S. 539 (1842)\par}
\vspace{0.8em}
{\large Case Brief\par}
\end{center}
\vspace{1em}

\section{Memory Jogger}
In \href{https://www.courtlistener.com/opinion/86231/prigg-v-pennsylvania/}{\textit{Prigg v. Pennsylvania}}, the Supreme Court struck down a Pennsylvania anti-kidnapping statute as violating the Fugitive Slave Clause and the Fugitive Slave Act of 1793, holding that Congress possesses exclusive legislative authority over fugitive-slave rendition---an expansive field-exclusivity model paralleling antebellum Commerce Clause doctrine---while simultaneously indicating that states cannot be compelled to assist in federal enforcement.

\section{Detailed Case Facts}
Margaret Morgan was held as a slave for life in Harford County, Maryland, under Maryland law, by Margaret Ashmore (\href{https://supreme.justia.com/cases/federal/us/41/539/}{Justia, 41 U.S. at 556, 608--609}). Although Morgan's former owner, John Ashmore, had allowed her to live in virtual freedom, she had never been formally emancipated (\href{https://www.oyez.org/cases/1789-1850/41us539}{Oyez}). In 1832, Morgan moved across the border into York County, Pennsylvania, without Margaret Ashmore's knowledge or consent (\href{https://supreme.justia.com/cases/federal/us/41/539/}{Justia, 41 U.S. at 556}). While residing in Pennsylvania, Morgan gave birth to at least one child more than a year after leaving Maryland (\href{https://supreme.justia.com/cases/federal/us/41/539/}{Justia, 41 U.S. at 557, 609}).

In February 1837, Margaret Ashmore formally appointed Edward Prigg, a Maryland citizen, as her agent and attorney to seize and return Margaret Morgan to Maryland (\href{https://supreme.justia.com/cases/federal/us/41/539/}{Justia, 41 U.S. at 556}). Prigg appeared before Thomas Henderson, a York County justice of the peace, swore out an oath that Morgan was a fugitive from labor, and obtained a warrant directing York County constable William McCleary to apprehend Morgan and her children (\href{https://supreme.justia.com/cases/federal/us/41/539/}{Justia, 41 U.S. at 556--557}). After McCleary brought Morgan and her children before Justice Henderson, Henderson refused to take further cognizance of the matter (\href{https://supreme.justia.com/cases/federal/us/41/539/}{Justia, 41 U.S. at 557}). Rather than pursuing judicial certification before a federal judge or complying with Pennsylvania's statutory procedures, Prigg---along with Nathan S. Bemis, Jacob Forward, and Stephen Lewis, Jr.---forcibly carried Morgan and her children out of Pennsylvania and delivered them to Ashmore in Maryland (\href{https://supreme.justia.com/cases/federal/us/41/539/}{Justia, 41 U.S. at 543, 557}).

At the time of the abduction, two legal frameworks governed the interstate recovery of alleged fugitive slaves. Nationally, Article IV, Section 2, Clause 3 of the U.S. Constitution (the Fugitive Slave Clause) provided that ``No Person held to Service or Labour in one State, under the Laws thereof, escaping into another, shall, in Consequence of any Law or Regulation therein, be discharged from such Service or Labour, but shall be delivered up on Claim of the Party to whom such Service or Labour may be due,'' and Congress had enacted the Fugitive Slave Act of February 12, 1793, ch. 7 (also cited in the opinion as ch. 51), 1 Stat. 302, empowering slaveholders or their agents to seize fugitives and bring them before any federal circuit or district judge or local state magistrate to obtain a certificate of removal (\href{https://supreme.justia.com/cases/federal/us/41/539/}{Justia, 41 U.S. at 611, 616--617}). At the state level, Pennsylvania had enacted its Act of March 25, 1826 (cited in Story's opinion as March 26, 1826), titled ``An act to give effect to the provisions of the Constitution of the United States relative to fugitives from labor, for the protection of free people of color, and prevent kidnapping,'' which made it a felony punishable by a fine of \$500 to \$1,000 and seven to twenty-one years of hard labor to forcibly or fraudulently remove any Black person from Pennsylvania with the intent to enslave them, and required claimants to obtain warrants and removal certificates through specified Pennsylvania judicial procedures while barring state aldermen and justices of the peace from exercising jurisdiction under the federal Act of 1793 (\href{https://supreme.justia.com/cases/federal/us/41/539/}{Justia, 41 U.S. at 550--555, 608}).

\section{Procedural History}
A York County grand jury indicted Edward Prigg and his three companions for kidnapping under Section 1 of Pennsylvania's 1826 statute (\href{https://supreme.justia.com/cases/federal/us/41/539/}{Justia, 41 U.S. at 543, 550}). To resolve growing interstate friction between Maryland and Pennsylvania, the two states cooperated to arrange a friendly test case: Pennsylvania enacted a special statute on May 22, 1839, and counsel for Maryland (representing Prigg) and the Attorney General of Pennsylvania agreed to submit a special verdict in the Court of Oyer and Terminer (also referred to in the record as the Court of Quarter Sessions) of York County (\href{https://supreme.justia.com/cases/federal/us/41/539/}{Justia, 41 U.S. at 543, 557--558, 609}). The jury returned the agreed special verdict finding the underlying facts of Morgan's escape, her Pennsylvania-born child, and Prigg's removal of Morgan without complying with the 1826 Act, and by agreement the trial court entered judgment of conviction against Prigg (\href{https://supreme.justia.com/cases/federal/us/41/539/}{Justia, 41 U.S. at 556--558}).

Prigg prosecuted a writ of error to the Supreme Court of Pennsylvania at the May Term of 1840, assigning as error that Pennsylvania's 1826 statute was repugnant to the U.S. Constitution (\href{https://supreme.justia.com/cases/federal/us/41/539/}{Justia, 41 U.S. at 558}). The Supreme Court of Pennsylvania affirmed the conviction \textit{pro forma} without opinion to expedite review (\href{https://supreme.justia.com/cases/federal/us/41/539/}{Justia, 41 U.S. at 558, 609}). Prigg then brought a writ of error to the U.S. Supreme Court under Section 25 of the Judiciary Act of 1789 (\href{https://supreme.justia.com/cases/federal/us/41/539/}{Justia, 41 U.S. at 539, 608}). The case was argued February 8--10, 1842, and decided March 1, 1842 (\href{https://supreme.justia.com/cases/federal/us/41/539/}{Justia, 41 U.S. at 539}; \href{https://www.oyez.org/cases/1789-1850/41us539}{Oyez}).

\section{Judicial Votes}
The Supreme Court voted \textbf{8--1} to reverse Edward Prigg's conviction and invalidate Pennsylvania's 1826 statute, producing seven separate opinions across the nine Justices (\href{https://www.courtlistener.com/opinion/86231/prigg-v-pennsylvania/}{CourtListener}; \href{https://blackpast.org/african-american-history/prigg-v-pennsylvania-1842/}{BlackPast}; \href{https://www.cambridge.org/core/journals/law-and-history-review/article/triumph-of-freedom-after-all-prigg-v-pennsylvania-reexamined/A0B7452002A367E4EFE865A3375F9D0E}{Goldstein, 29 Law \& Hist. Rev. 763 (2011)}):
\begin{itemize}
\item \textbf{Opinion of the Court:} Justice Joseph Story delivered the opinion of the Court, joined silently in full by Justice John Catron and Justice John McKinley, and joined in full by Justice James Moore Wayne (who also wrote separately).
\item \textbf{Concurring:} Justice James Moore Wayne concurred ``altogether'' in Story's opinion and wrote separately to reinforce federal legislative exclusivity.
\item \textbf{Concurring in Part and Dissenting in Part (Concurring in Judgment):} Chief Justice Roger B. Taney concurred in reversing the conviction and invalidating the Pennsylvania statute, but dissented from Story's holding that Congress's power is exclusive and prohibits state laws aiding slaveholders.
\item \textbf{Concurring in Judgment:} Justice Smith Thompson concurred in the judgment because the 1826 statute conflicted with the federal Act of 1793, but rejected Story's doctrine of exclusive federal legislative power in the absence of congressional legislation.
\item \textbf{Concurring in Judgment:} Justice Henry Baldwin concurred in reversing the conviction on the ground that removing an admitted slave could not be kidnapping, but dissented from the Court's reasoning, viewing the Fugitive Slave Clause as completely self-executing without needing congressional legislation (though agreeing that if legislation were necessary, it would belong exclusively to Congress).
\item \textbf{Concurring in Part and Dissenting in Part (Concurring in Judgment):} Justice Peter V. Daniel concurred in invalidating Pennsylvania's statute and reversing the conviction, but dissented from the exclusivity doctrine, arguing that states retain concurrent or ancillary power to legislate in aid of slaveholders.
\item \textbf{Dissenting:} Justice John McLean was the sole dissenter from the judgment of reversal, agreeing that federal legislative power is exclusive but rejecting the constitutional right of forcible private recaption without judicial process and defending Pennsylvania's anti-kidnapping statute as a valid state police regulation.
\end{itemize}
On the pivotal federalism question of whether congressional power to legislate over fugitive-slave rendition is \textit{exclusive} of state legislation, six Justices (Story, Catron, McKinley, Wayne, McLean, and conditionally Baldwin) supported exclusivity, whereas three (Taney, Thompson, and Daniel) insisted states retained concurrent authority to enact legislation assisting slaveholders (\href{https://www.cambridge.org/core/journals/law-and-history-review/article/triumph-of-freedom-after-all-prigg-v-pennsylvania-reexamined/A0B7452002A367E4EFE865A3375F9D0E}{Goldstein, 29 Law \& Hist. Rev. 763 (2011)}).

\section{Holding}
\textbf{Judgment:} Reversed and remanded to the Supreme Court of Pennsylvania with directions to enter judgment of ``not guilty'' in favor of Edward Prigg on the special verdict (\href{https://supreme.justia.com/cases/federal/us/41/539/}{Justia, 41 U.S. at 626}; \href{https://blackpast.org/african-american-history/prigg-v-pennsylvania-1842/}{BlackPast}).

\textbf{Rule Announced:} First, Article IV, Section 2, Clause 3 confers a positive, unqualified, self-executing constitutional right upon slaveholders to peaceably seize and recapture fugitive slaves in any state without illegal violence or breach of the peace, rendering void any state law that interrupts, delays, or punishes that recaption (\href{https://supreme.justia.com/cases/federal/us/41/539/}{Justia, 41 U.S. at 612--613, 625--626}). Second, Congress has implied constitutional authority to enact the Fugitive Slave Act of 1793, and that legislative power is \textit{exclusive} in the National Government, completely preempting state legislation on the subject of fugitive-slave rendition---much like exclusive national powers recognized under Commerce Clause and bankruptcy precedents such as \textit{Sturgis v. Crowninshield} and \textit{Houston v. Moore} (\href{https://supreme.justia.com/cases/federal/us/41/539/}{Justia, 41 U.S. at 617--618, 622--625}). Third, states may not be compelled to enforce the federal Fugitive Slave Act, and although state magistrates may voluntarily exercise authority conferred by the 1793 Act, state legislatures may prohibit their own officials from doing so (\href{https://supreme.justia.com/cases/federal/us/41/539/}{Justia, 41 U.S. at 615--616, 622}).

\textbf{Holding vs. Dicta:} The strict holding necessary to reverse Prigg's conviction was that Pennsylvania's 1826 statute was unconstitutional insofar as it criminalized a slaveholder's agent for removing an admitted fugitive slave in conflict with the Fugitive Slave Clause and the 1793 Act. As Chief Justice Taney and Justice Daniel expressly objected, Story's broader pronouncements---that federal legislative power is completely \textit{exclusive} even in the absence of congressional action, that states may not enact laws \textit{aiding} slaveholders, that states retain general ``police power'' to arrest and expel runaways as vagrants, and that state legislatures may forbid state magistrates from enforcing the 1793 Act---went beyond the narrow conflict before the Court (\href{https://supreme.justia.com/cases/federal/us/41/539/}{Justia, 41 U.S. at 622, 625, 627}; \href{https://blackpast.org/african-american-history/prigg-v-pennsylvania-1842/}{BlackPast}). Nevertheless, Story's exclusivity rationale and his non-compulsion language regarding state magistrates proved enormously consequential, directly spurring Northern ``personal liberty laws'' withdrawing state enforcement assistance and later shaping the Fugitive Slave Act of 1850 (\href{https://ballotpedia.org/Prigg_v._Pennsylvania}{Ballotpedia}; \href{https://www.encyclopedia.com/history/dictionaries-thesauruses-pictures-and-press-releases/prigg-v-commonwealth-pennsylvania}{Encyclopedia.com}).

\section{Analysis of Opinions}

\subsection{Majority Opinion (Justice Story)}
Justice Story began by framing the Fugitive Slave Clause as a ``fundamental article without the adoption of which the Union could not have been formed'' (\href{https://supreme.justia.com/cases/federal/us/41/539/}{Justia, 41 U.S. at 611}). Invoking Lord Mansfield's decision in \textit{Somerset's Case}, Story acknowledged that under the general law of nations, slavery is ``a mere municipal regulation founded upon and limited to the range of the territorial laws,'' meaning free states would otherwise be under no obligation to recognize slave status (\href{https://supreme.justia.com/cases/federal/us/41/539/}{Justia, 41 U.S. at 611--612}). Because the Constitution displaced comity with a ``positive, unqualified right,'' any state law that ``interrupts, limits, delays, or postpones'' a master's immediate possession operates \textit{pro tanto} as a forbidden discharge (\href{https://supreme.justia.com/cases/federal/us/41/539/}{Justia, 41 U.S. at 612--613}). Drawing on Blackstone's common-law doctrine of recaption, Story concluded that the Clause is self-executing to the extent that a master may peaceably seize and recapture a fugitive anywhere in the Union without a breach of the peace (\href{https://supreme.justia.com/cases/federal/us/41/539/}{Justia, 41 U.S. at 613}).

Turning to the Clause's command that a fugitive ``shall be delivered up on claim,'' Story reasoned that because ``where the end is required, the means are given,'' Congress possesses implied legislative power to enact enforcement remedies, confirmed by the 1793 Act and decades of judicial acquiescence (\href{https://supreme.justia.com/cases/federal/us/41/539/}{Justia, 41 U.S. at 615--622}). Crucially for federalism and Commerce Clause jurisprudence, Story then held that this federal legislative power is \textit{exclusive}. Relying on \textit{Houston v. Moore}, 18 U.S. (5 Wheat.) 1, 21--22 (1820), and Chief Justice Marshall's exclusivity test in \textit{Sturgis v. Crowninshield}, 17 U.S. (4 Wheat.) 122, 193 (1819), Story argued that when Congress regulates a subject, its silence is as expressive as its enactments, prohibiting states from adding ``auxiliary'' regulations (\href{https://supreme.justia.com/cases/federal/us/41/539/}{Justia, 41 U.S. at 617--618, 622--625}). Because the right derived solely from the national Constitution and required national uniformity across state lines, state concurrent authority would breed ``endless'' confusion (\href{https://supreme.justia.com/cases/federal/us/41/539/}{Justia, 41 U.S. at 622--625}).

At the same time, Story introduced two vital limits. First, because the Clause appears in the national Constitution and assigns no duties to state officers, states ``cannot, therefore, be compelled to enforce them,'' and state magistrates may act under the 1793 Act only ``unless prohibited by state legislation'' (\href{https://supreme.justia.com/cases/federal/us/41/539/}{Justia, 41 U.S. at 615--616, 622}). Second, Story reserved the states' ``police power'' to arrest, restrain, and remove runaway slaves to guard against ``depredations and evil example,'' provided such measures do not obstruct the owner's federal right (\href{https://supreme.justia.com/cases/federal/us/41/539/}{Justia, 41 U.S. at 625}). Because Pennsylvania's 1826 statute punished the very act of recaption authorized by the Constitution, it was void (\href{https://supreme.justia.com/cases/federal/us/41/539/}{Justia, 41 U.S. at 625--626}). Notably, Story's opinion ignored the special verdict's finding that one of Morgan's children had been born free on Pennsylvania soil.

\subsection{Concurring Opinion (Justice Wayne)}
Justice Wayne concurred fully in Story's opinion and enumerated seven propositions settled by the Court (\href{https://supreme.justia.com/cases/federal/us/41/539/}{Justia, 41 U.S. at 636--637}). Writing separately to defend federal exclusivity, Wayne invoked the law of nations governing treaties and interstate compacts: because the Fugitive Slave Clause was an indispensable constitutional compromise, neither side to the bargain could unilaterally interpret or modify its execution through state remedies (\href{https://supreme.justia.com/cases/federal/us/41/539/}{Justia, 41 U.S. at 638--643}). Even seemingly helpful state remedies, Wayne warned, carried the power to impose delays, expenses, and hostile local procedures that would effectively destroy slaveholders' property rights (\href{https://supreme.justia.com/cases/federal/us/41/539/}{Justia, 41 U.S. at 643--648}).

\subsection{Concurring in Part and Dissenting in Part (Chief Justice Taney)}
Chief Justice Taney agreed that the master possessed a constitutional right of peaceable recaption, that the 1793 Act was valid, and that Pennsylvania's 1826 law was unconstitutional (\href{https://supreme.justia.com/cases/federal/us/41/539/}{Justia, 41 U.S. at 626--627}). However, Taney sharply dissented from Story's exclusivity holding. Textually, the Clause forbids only state laws that \textit{impair} or discharge the master's right; it contains no prohibition against state laws that \textit{support} and enforce it (\href{https://supreme.justia.com/cases/federal/us/41/539/}{Justia, 41 U.S. at 627--629}). Analogizing to the Contracts Clause and the Privileges and Immunities Clause, Taney argued that states have not only the power but an affirmative constitutional duty to protect federally guaranteed rights (\href{https://supreme.justia.com/cases/federal/us/41/539/}{Justia, 41 U.S. at 628--629}). Practically, Taney warned that because there were very few federal judges in the nation, barring state officers and state laws from assisting slaveholders would render the 1793 Act ``ineffectual and delusive'' and turn border states like Maryland into ``an open pathway for the fugitives escaping from other states'' (\href{https://supreme.justia.com/cases/federal/us/41/539/}{Justia, 41 U.S. at 630--632}). He also rejected Story's police-power workaround as untenable when a fugitive was traveling peaceably through a state (\href{https://supreme.justia.com/cases/federal/us/41/539/}{Justia, 41 U.S. at 632--633}).

\subsection{Concurring in Judgment (Justice Thompson)}
Justice Thompson concurred in reversing Prigg's conviction because Pennsylvania's statute conflicted with the 1793 Act, which prevailed under the Supremacy Clause (\href{https://supreme.justia.com/cases/federal/us/41/539/}{Justia, 41 U.S. at 634--635}). Like Taney, however, Thompson rejected federal exclusivity in the absence of congressional legislation. He argued that nothing in the text or nature of the Fugitive Slave Clause barred states from passing laws in affirmance of the master's right if Congress had not acted or were to repeal the 1793 Act (\href{https://supreme.justia.com/cases/federal/us/41/539/}{Justia, 41 U.S. at 635--636}).

\subsection{Concurring in Judgment (Justice Baldwin)}
In a brief reporter's summary, Justice Baldwin concurred in reversing the conviction on the ground that because Margaret Morgan's status as a slave was admitted on the record, removing her could not constitute kidnapping (\href{https://supreme.justia.com/cases/federal/us/41/539/}{Justia, 41 U.S. at 636}). Baldwin dissented from the Court's broader principles, maintaining (as Justice Wayne summarized) that the Fugitive Slave Clause was completely self-executing and conferred no legislative power on Congress, though if congressional legislation were necessary, that power would be exclusive (\href{https://supreme.justia.com/cases/federal/us/41/539/}{Justia, 41 U.S. at 636--637}).

\subsection{Concurring in Part and Dissenting in Part (Justice Daniel)}
Justice Daniel concurred in invalidating Pennsylvania's statute as conflicting with the Constitution and the 1793 Act, but dissented from Story's exclusivity doctrine (\href{https://blackpast.org/african-american-history/prigg-v-pennsylvania-1842/}{BlackPast, 41 U.S. at 650--658}). Invoking Commerce Clause and bankruptcy federalism precedents---including \textit{Sturgis v. Crowninshield}, Story's own dissent in \textit{Houston v. Moore}, \textit{Gibbons v. Ogden}, \textit{Brown v. Maryland}, and \textit{City of New York v. Miln}---Daniel argued that mere grants of federal power do not strip states of concurrent or ancillary authority unless the Constitution expressly makes the power exclusive or state action directly collides with federal law (\href{https://blackpast.org/african-american-history/prigg-v-pennsylvania-1842/}{BlackPast, 41 U.S. at 652--656}). Like Taney, Daniel warned that Story's exclusivity rule would neutralize state assistance and that invoking the police power to deport runaways would push fugitives farther from their masters rather than securing their return (\href{https://blackpast.org/african-american-history/prigg-v-pennsylvania-1842/}{BlackPast, 41 U.S. at 656--658}).

\subsection{Dissenting Opinion (Justice McLean)}
Justice McLean alone dissented from the judgment reversing Prigg's conviction (\href{https://blackpast.org/african-american-history/prigg-v-pennsylvania-1842/}{BlackPast, 41 U.S. at 658--674}). Notably, McLean \textit{agreed} with Story that Congress's legislative power over fugitive-slave rendition is exclusive---analogizing it directly to ``the power to regulate commerce, or that of foreign intercourse''---and went even further than Story by arguing that Congress \textit{could} impose mandatory duties on state magistrates under the 1793 Act (\href{https://blackpast.org/african-american-history/prigg-v-pennsylvania-1842/}{BlackPast, 41 U.S. at 660--666}). Where McLean broke decisively with the majority was on the right of private forcible recaption. Reading the 1793 Act to cover the entire field, McLean insisted that a slaveholder who seizes an alleged fugitive must bring that person before a federal judge or magistrate to obtain a certificate before removing them from the state (\href{https://blackpast.org/african-american-history/prigg-v-pennsylvania-1842/}{BlackPast, 41 U.S. at 666--671}). Because a free state presumes every resident within its borders to be free regardless of color, Pennsylvania's anti-kidnapping statute was a valid police regulation protecting free Black citizens against unauthorized abduction; when the local magistrate refused to act, Prigg was required to take Morgan before one of the two federal judges residing in Pennsylvania rather than abducting her by force (\href{https://blackpast.org/african-american-history/prigg-v-pennsylvania-1842/}{BlackPast, 41 U.S. at 669--674}).

\section{Examples: Future Applications}
Although the Thirteenth Amendment abolished slavery and rendered the Fugitive Slave Clause obsolete, \textit{Prigg}'s structural doctrines---field exclusivity, preemption of ``auxiliary'' state regulations, and non-compulsion of state officials---map directly onto modern Commerce Clause and federal preemption disputes:
\begin{enumerate}
\item \textbf{Same-Side (Invalidated under Exclusivity/Preemption):} Suppose Congress enacts a comprehensive federal statute regulating interstate railroad safety standards and locomotive equipment specifications under the Commerce Clause. A state passes a supplementary law imposing additional inspection certificates and waiting periods before trains may cross its borders. Under \textit{Prigg}'s exclusivity and field-preemption logic (echoing \textit{Houston v. Moore}), the state's ``auxiliary'' regulation is invalid because Congress's silence on additional requirements is as expressive as its text, and state-by-state variations destroy national uniformity.
\item \textbf{Same-Side (Invalidated under Exclusivity/Preemption):} Suppose a state enacts a criminal statute prosecuting federal immigration enforcement agents or private contractors who transport detainees across state lines pursuant to federal immigration warrants unless they first obtain a state magistrate's removal certificate. Under \textit{Prigg}, because immigration and interstate rendition under federal authority belong exclusively to the national government, the state cannot criminalize or delay the exercise of federally authorized removal by interposing state procedural hurdles.
\item \textbf{Opposite-Side (Upheld under State Police Power / Non-Compulsion):} Suppose Congress passes a federal statute authorizing state and local police officers to enforce federal immigration or commercial-licensing rules, and a state legislature responds by enacting a law prohibiting its own state and municipal officers from expending state resources to assist in enforcing that federal program (without obstructing federal agents themselves). Under \textit{Prigg}'s non-compulsion principle---``unless prohibited by state legislation''---the state law is valid because Congress cannot compel state officials to execute federal duties, and a state's withdrawal of voluntary assistance does not affirmatively obstruct federal officers.
\item \textbf{Opposite-Side (Upheld under State Police Power / Concurrent Authority):} Suppose a state enacts a general highway-safety and anti-breach-of-the-peace ordinance prohibiting any person from using reckless force or blocking public thoroughfares while repossessing commercial vehicles engaged in interstate commerce. Under \textit{Prigg}'s express reservation of the state police power and its limitation of private self-help to recaption ``without any breach of the peace or any illegal violence,'' the neutral safety law is valid because it protects local public order without targeting or discriminating against the federal right.
\item \textbf{Fence-Sitter (Borderline / Contested):} Suppose Congress regulates interstate transportation of hazardous materials under the Commerce Clause without expressly preempting state law, and a state passes a ``cooperative'' statute that imposes identical substantive standards as federal law while adding a state-court civil penalty and mandatory state-level recordkeeping to ``aid'' federal enforcement. Under Story's strict exclusivity framework in \textit{Prigg}, even well-intentioned state legislation in ``aid'' of an exclusive federal domain is void because Congress's chosen remedies occupy the field; under Taney's and Daniel's concurrent-power approach (and modern Commerce Clause preemption doctrine requiring clearer congressional intent to occupy a field), the state's complementary enforcement might survive so long as it does not conflict with federal objectives.
\end{enumerate}

\section{Critique}
\textit{Prigg v. Pennsylvania} has provoked enduring debate across originalist, textualist, and progressive-doctrinal scholarship (\href{https://www.cambridge.org/core/journals/law-and-history-review/article/triumph-of-freedom-after-all-prigg-v-pennsylvania-reexamined/A0B7452002A367E4EFE865A3375F9D0E}{Goldstein, 29 Law \& Hist. Rev. 763 (2011)}).

\textbf{Originalist and Textualist Perspectives:} Textualist and originalist critics note that Article IV, Section 2, Clause 3 contains no express grant of legislative power to Congress (unlike the Full Faith and Credit Clause in Article IV, Section 1) and no express prohibition on state laws aiding rendition (\href{https://chicagounbound.uchicago.edu/cgi/viewcontent.cgi?article=6111&context=journal_articles}{Currie, \textit{Article IV and Federal Powers}}; \href{https://scholarship.law.wm.edu/cgi/viewcontent.cgi?article=2141&context=wmborj}{Barnett \& Blackman, Wm. \& Mary Bill Rts. J.}). Chief Justice Taney and Justice Daniel pointed out that Story inferred both an unenumerated congressional power and its exclusivity from textual silence. Historians such as Don E. Fehrenbacher have also questioned Story's sweeping historical assertion that the Union could never have been formed without the Fugitive Slave Clause, noting that the Clause was added late in the Constitutional Convention with minimal debate (\href{https://en.wikipedia.org/wiki/Prigg_v._Pennsylvania}{Wikipedia, citing Fehrenbacher}). From a Commerce Clause doctrinal standpoint, legal scholar David Currie argued that ``[n]one of Story's three arguments will bear close scrutiny,'' observing that Pennsylvania's 1826 statute did not discharge slaves but merely established orderly procedures to verify claims (\href{https://chicagounbound.uchicago.edu/cgi/viewcontent.cgi?article=6111&context=journal_articles}{Currie}).

\textbf{Progressive and Historical-Doctrinal Perspectives:} Revisionist legal historians, led by Paul Finkelman and Barbara Holden-Smith, sharply reject William Wetmore Story's nineteenth-century defense that his father's opinion was a covert ``triumph of freedom'' designed to hobble slave catching (\href{https://chicagounbound.uchicago.edu/supremecourtrev/vol1994/iss1/7/}{Finkelman, 1994 Sup. Ct. Rev. 247}; \href{https://scholarship.law.cornell.edu/clr/vol78/iss6/2/}{Holden-Smith, 78 Cornell L. Rev. 1086 (1993)}). Finkelman demonstrates that Story prioritized judicial nationalism and pro-slavery federal supremacy over the civil liberties of free Black Northerners, stripping Northern states of personal liberty laws that prevented kidnapping and even privately writing to Senator John Macpherson Berrien eight weeks after \textit{Prigg} to recommend creating federal commissioners to facilitate slave recapture---a recommendation Congress later adopted in the draconian Fugitive Slave Act of 1850 (\href{https://msa.maryland.gov/megafile/msa/speccol/sc5400/sc5496/051200/051268/images/finkelman_story.pdf}{Finkelman, 1994 Sup. Ct. Rev. 247}). By contrast, scholars such as Leslie Friedman Goldstein argue that Story's adoption of \textit{Somerset}'s principle that slavery is merely local ``municipal'' law---combined with his invitation for Northern states to bar their officials from enforcing federal slave laws---genuinely bolstered the antislavery constitutionalism of ``freedom national, slavery local'' (\href{https://www.cambridge.org/core/journals/law-and-history-review/article/triumph-of-freedom-after-all-prigg-v-pennsylvania-reexamined/A0B7452002A367E4EFE865A3375F9D0E}{Goldstein, 29 Law \& Hist. Rev. 763 (2011)}).

\textbf{Competing Underlying Values:} \textit{Prigg} exposed an irreconcilable collision among three constitutional values: (1) \textit{national uniformity and federal supremacy}, which Story elevated through an aggressive field-exclusivity model akin to Dormant Commerce Clause exclusivity; (2) \textit{state sovereignty and police-power protection of individual liberty}, which Justice McLean championed to shield free Black citizens from summary abduction; and (3) \textit{sectional compromise over chattel slavery}, which ultimately proved incapable of judicial resolution.

\section{Key Quotations}

\begin{quote}
``The clause manifestly contemplates the existence of a positive, unqualified right on the part of the owner of the slave which no state law or regulation can in any way qualify, regulate, control, or restrain. The slave is not to be discharged from service or labor in consequence of any state law or regulation. Now certainly, without indulging in any nicety of criticism upon words, it may fairly and reasonably be said that any state law or state regulation which interrupts, limits, delays, or postpones the right of the owner to the immediate possession of the slave and the immediate command of his service and labor operates pro tanto a discharge of the slave therefrom.''
\end{quote}
\noindent --- Justice Joseph Story, Opinion of the Court, 41 U.S. at 612--613 (\href{https://supreme.justia.com/cases/federal/us/41/539/}{Justia}; \href{https://www.courtlistener.com/opinion/86231/prigg-v-pennsylvania/}{CourtListener}).

\begin{quote}
``Upon this ground, we have not the slightest hesitation in holding that, under and in virtue of the Constitution, the owner of a slave is clothed with entire authority, in every State in the Union, to seize and recapture his slave whenever he can do it without any breach of the peace or any illegal violence. In this sense and to this extent, this clause of the Constitution may properly be said to execute itself, and to require no aid from legislation, state or national.''
\end{quote}
\noindent --- Justice Joseph Story, Opinion of the Court, 41 U.S. at 613 (\href{https://supreme.justia.com/cases/federal/us/41/539/}{Justia}; \href{https://www.courtlistener.com/opinion/86231/prigg-v-pennsylvania/}{CourtListener}).

\begin{quote}
``As to the authority so conferred upon state magistrates, while a difference of opinion has existed, and may exist still, on the point in different States, whether state magistrates are bound to act under it, none is entertained by this Court that state magistrates may, if they choose, exercise that authority unless prohibited by state legislation.''
\end{quote}
\noindent --- Justice Joseph Story, Opinion of the Court, 41 U.S. at 622 (\href{https://supreme.justia.com/cases/federal/us/41/539/}{Justia}; \href{https://www.courtlistener.com/opinion/86231/prigg-v-pennsylvania/}{CourtListener}).

\begin{quote}
``The language used in the Constitution does not, in my judgment, justify this construction given to it by the court. It contains no words prohibiting the several States from passing laws to enforce this right. They are, in express terms, forbidden to make any regulation that shall impair it, but there the prohibition stops.''
\end{quote}
\noindent --- Chief Justice Roger B. Taney, Concurring in Part and Dissenting in Part, 41 U.S. at 627--628 (\href{https://supreme.justia.com/cases/federal/us/41/539/}{Justia}; \href{https://www.courtlistener.com/opinion/86231/prigg-v-pennsylvania/}{CourtListener}).

\begin{quote}
``In a state where slavery is allowed, every coloured person is presumed to be a slave; and on the same principle, in a non-slaveholding state, every person is presumed to be free without regard to colour.''
\end{quote}
\noindent --- Justice John McLean, Dissenting, 41 U.S. at 669 (\href{https://blackpast.org/african-american-history/prigg-v-pennsylvania-1842/}{BlackPast}; \href{https://www.courtlistener.com/opinion/86231/prigg-v-pennsylvania/}{CourtListener}).

\end{document}
